% STATUS: rewritten 2026-11-05
% STATUS: extended 2026-11-07 (Fosner and Fosner 2014, Pandey 2024 and
% Di Bartolo, Di Fatta and La Rosa 2025 superbiderivations read in full)
% Sources read in full: Musson 2012 (Definition of a Lie superalgebra, p. 1;
% Lemma 1.1.2, pp. 2--3; Theorem 2.6.2, pp. 19--20), Kac 1977 (Theorem 1,
% pp. 1385--1386); Chi--Zhang--Meng--Tang 2026 (Theorems 3.2, 3.5 and 4.2,
% pp. 7, 12, 15; read 2026-11-06); Fosner and Fosner 2014 (pp. 303--306);
% Pandey 2024 (pp. 1--2, 7--8); Di Bartolo, Di Fatta and La Rosa 2025
% superbiderivations (pp. 1--2, 8--10).
% Sources: Musson2012, Kac1977Classification, ChiZhangMengTang2026.
\section{Lie and Jordan superalgebras}\label{sec:nonassociative:superalgebras}

The Leibniz algebras of the previous section keep one product rule. % claim-ok: own summary of the Leibniz section, links to Musson below
A superalgebra separates the even part from the odd part, and the
product rule picks up a sign \cite[pp.~1--2]{Musson2012}.

Each result below concerns derivations of a graded algebra, where the product rule picks up the sign of the two degrees involved.

\begin{definition}[{\cite[pp.~1--2]{Musson2012}}]
\label{def:nonassociative:superderivation}
Let $A$ be a $\mathbb{Z}_2$-graded algebra over a field $K$ that need
not be associative. A $K$-linear map $\partial\colon A\to A$ of degree
$\alpha\in\mathbb{Z}_2$ is a \emph{left superderivation} if
\[
\partial(bc)\;=\;\partial(b)c+(-1)^{\alpha\bar{b}}\,b\,\partial(c)
\qquad\text{for all }b,c\in A,
\]
where $\bar b$ is the degree of the homogeneous element $b$. A
\emph{superderivation} is a left superderivation. A superderivation of
degree $0$ is called \emph{even}, and one of degree $1$ is called
\emph{odd}.
\end{definition}

The graded bracket of the next definition supplies the maps on which
superderivations act. % claim-ok: own linking sentence

\begin{definition}[{\cite[p.~1]{Musson2012}}]
\label{def:nonassociative:lie-super}
A \emph{Lie superalgebra} is a $\mathbb{Z}_2$-graded vector space
$\mathfrak{g}=\mathfrak{g}_0\oplus\mathfrak{g}_1$ over $K$ with a
bilinear map $[-,-]\colon\mathfrak{g}\times\mathfrak{g}\to\mathfrak{g}$
that is graded skew-symmetric and satisfies the graded Jacobi identity.
\end{definition}

Musson proves the following equivalence (Lemma~\ref{lem:nonassociative:musson-jacobi}) \cite[Lemma~1.1.2,
pp.~2--3]{Musson2012}.\begin{lemma}[{\cite[Lemma~1.1.2, pp.~2--3]{Musson2012}}]
\label{lem:nonassociative:musson-jacobi}
Let $A$ be a $\mathbb{Z}_2$-graded algebra whose product
$(a,b)\mapsto [a,b]$ is graded skew-symmetric. The following conditions
are equivalent.
\begin{enumerate}
\item $A$ satisfies the graded Jacobi identity.
\item For every $a\in A$ the map $\operatorname{ad} a: b\mapsto [a,b]$
is a left superderivation of degree $\bar a$.
\item For every homogeneous $a\in A_{\alpha}$ the map
$b\mapsto [b,a]$ is a right superderivation of degree $\alpha$.
\end{enumerate}
\end{lemma}

Musson quotes the following theorem of Scheunert, where a nondegenerate
Killing form stands in place of simplicity (Theorem~\ref{thm:nonassociative:super-inner})
\cite[Theorem~2.6.2, pp.~19--20]{Musson2012}.\begin{theorem}[{\cite[Theorem~2.6.2, pp.~19--20]{Musson2012}}]
\label{thm:nonassociative:super-inner}
Let $\mathfrak{g}$ be a Lie superalgebra whose Killing form is
nondegenerate. Then every superderivation of $\mathfrak{g}$ is inner,
that is, of the form $\operatorname{ad} a$ for some $a\in\mathfrak{g}$.
\end{theorem}

Musson points out that the superalgebras with a nondegenerate Killing
form behave like semisimple Lie algebras
\cite[p.~20]{Musson2012}. Kac classified the simple finite-dimensional
Lie superalgebras as follows (Theorem~\ref{thm:nonassociative:kac-classification}) \cite[Theorem~1,
pp.~1385--1386]{Kac1977Classification}.

The classification enters this chapter through one observation: a degree-zero
derivation of a graded Lie superalgebra is an infinitesimal automorphism of the
grading. Listing the gradings therefore lists the derivations up to the
automorphism group, which is what the superbiderivations below are computed
against.
\begin{theorem}[{\cite[Theorem~1,
pp.~1385--1386]{Kac1977Classification}}]
\label{thm:nonassociative:kac-classification}
For a simple finite-dimensional Lie superalgebra over an algebraically
closed field of characteristic $0$, all $\mathbb{Z}$-gradings are known
explicitly: the list in \cite{Kac1977Classification} includes every
$\mathbb{Z}$-grading (up to automorphism) of every such superalgebra.
\end{theorem}

A local superderivation matches its value at one point at a time.
Chi, Zhang, Meng and Tang work over a field of characteristic not two with
a finite-dimensional nilpotent Lie superalgebra
\cite[p.~1]{ChiZhangMengTang2026}. A local superderivation is called pure
when it is not a superderivation (Theorem~\ref{thm:nonassociative:pure-local-super}) \cite[p.~1]{ChiZhangMengTang2026}.

\begin{theorem}[{\cite[Theorems~3.2 and 3.5, pp.~7, 12]{ChiZhangMengTang2026}}]\label{thm:nonassociative:pure-local-super}
Let $L$ be a finite-dimensional two-step nilpotent Lie superalgebra over a
field of characteristic not two. Then $L$ admits a pure local
superderivation and a pure local anti-superderivation.
\end{theorem}

So on the nilpotent side local maps are genuinely wider than global ones.
The authors give sufficient conditions for the same failure at higher
nilpotency index \cite[p.~1]{ChiZhangMengTang2026}. Every three-step
nilpotent Lie superalgebra admits a pure local superderivation
\cite[Theorem~4.2, p.~15]{ChiZhangMengTang2026}.

A superderivation of a prime superalgebra obeys a polynomial identity only in
three ways. A superderivation of a superalgebra keeps the grading, so it is a
sum $D=D_0+D_1$ of an even and an odd part \cite[p.~304]{FosnerFosner2014}.
Fošner and Fošner impose two commutator identities on each part (Theorem~\ref{thm:nonassociative:fosner-super}).

\begin{theorem}[{\cite[Theorem~1, p.~305]{FosnerFosner2014}}]\label{thm:nonassociative:fosner-super}
Let $A$ be a prime associative superalgebra and let $D=D_0+D_1$ be a nonzero
superderivation of $A$. Suppose that
$D_i(x)[D_i(x),x]_s=0$ and $[D_i(x),x]_sD_i(x)=0$ for every homogeneous
element $x$ and $i=0,1$. Then one of the following holds:
\begin{enumerate}
\item $A_1=0$ and $D(A_0)\subseteq Z(A)$;
\item $D(A_0)=0$ and $A$ is commutative;
\item $D^{2}=0$.
\end{enumerate}
\end{theorem}

Here $[a,b]_s=ab-(-1)^{\bar a\bar b}ba$ is the super commutator
\cite[p.~304]{FosnerFosner2014}. The paper proves the analogous statements for
each degree separately, and the theorem is their union
\cite[Propositions~1--2, pp.~305--306]{FosnerFosner2014}.

The bilinear form of the question has its own definition for superalgebras. A
super-biderivation is a bilinear map that is a superderivation in each argument,
with the super sign carried through \cite[p.~2]{Pandey2024Superbiderivations}.
On a perfect Hom-Lie superalgebra it satisfies the relation
$\varphi(\omega,[\zeta,\xi])=(-1)^{|\varphi||\omega|}[\omega,\varphi(\zeta,\xi)]$
\cite[Lemma~3.7, p.~7]{Pandey2024Superbiderivations}. On a centerless algebra it
has the centroid form (Theorem~\ref{thm:nonassociative:hom-lie-superbiderivation}).

\begin{theorem}[{\cite[Theorem~3.8, p.~7]{Pandey2024Superbiderivations}}]\label{thm:nonassociative:hom-lie-superbiderivation}
Let $H$ be a centerless Hom-Lie superalgebra with $H=H'$. Then every
super-biderivation $\varphi\colon H\times H\to H$ has the form
$\varphi(\zeta,\xi)=\alpha^{-1}\delta([\zeta,\xi])$ for a map $\delta$ in the
centroid of $H$.
\end{theorem}

So the Hom-Lie superalgebra case has the shape of the Lie algebra case in
Theorem~\ref{thm:lie:tang-biderivation}, with the twist $\alpha^{-1}$ and a
centroid element in place of a scalar
\cite[Theorem~3.8, p.~7]{Pandey2024Superbiderivations}. A special
super-biderivation of a perfect centerless algebra is zero
\cite[Theorem~3.10, p.~9]{Pandey2024Superbiderivations}.

The definition of a superbiderivation itself has more than one form in the
literature. Di Bartolo, Di Fatta and La Rosa propose one that requires the
bracket of the superalgebra to be a superbiderivation, a condition the earlier
definitions do not meet \cite[p.~2]{DiBartoloDiFattaLaRosa2025Super}. Their two
equivalent forms are Definitions~2.5 and~2.6
\cite[Definitions~2.5--2.6, p.~5]{DiBartoloDiFattaLaRosa2025Super}. A nonzero bilinear map of odd
degree cannot be a superbiderivation under both the new and the earlier
definitions at once \cite[p.~7]{DiBartoloDiFattaLaRosa2025Super}. On a complete
Lie superalgebra every superbiderivation has the standard shape
\cite[Theorem~3.6, p.~10]{DiBartoloDiFattaLaRosa2025Super}.

With the classification in place we return to a map that appears both in
the alternative picture and in the Jordan superalgebra picture: the
$\delta$-derivation of Section~\ref{sec:nonassociative:delta-quasi-anti}.

Han proves that the derivations of a rational vertex operator algebra are
inner, the vertex-algebra analogue of the same question for algebras
\cite[Abstract, p.~1]{Han2026VOA}.

A graded ring gives both an ordinary and a graded answer. Ghahramani, Ghosseiri
and Safari characterize the derivations and biderivations of the quaternion ring
$H(S)$ over a ring $S$ in which 2 is invertible, and, reading $H(S)$ as a
$\mathbb{Z}_2$-graded ring, its superbiderivations
\cite[Abstract]{GhahramaniGhosseiriSafari2020}.

For the orthosymplectic Lie superalgebra $\mathfrak{osp}(1,2)$ over an algebraically closed field of characteristic $p>2$, Li, Liu and Wang give a necessary and sufficient condition for a map of the cartesian product of a finite-dimensional Lie superalgebra with itself into a nontrivial simple module to be a symmetric super-biderivation, and determine all super-biderivations from $\mathfrak{osp}(1,2)$ to its finite-dimensional simple modules~\cite{WlSuperBiderivationsFromOsp12ToA2025}.

Cheng and Sun study the super-biderivations of the twisted $N=2$ superconformal algebras, the notion being the Lie-superalgebra counterpart of biderivations of Lie algebras, and determine the linear super-commuting maps on them~\cite{WlSuperBiderivationsAndLinearSup2019}.

Cheraghpour and Ghosseiri describe the superderivations of trivial extensions and of upper triangular matrix rings, answering a question of Ghahramani and others in that setting, and show that every super-biderivation of these rings is a biderivation~\cite{WlOnSuperderivationsAndSuperBide2019}.

For a field of characteristic $p>3$, Ma and Chen study the derivations of the even part $\overline{M}$ of the finite-dimensional simple modular Lie superalgebra $M$ and determine its generator sets~\cite{WlDerivationsOfTheEvenPartOfFini2015}.

The outer derivation algebras of the finite-dimensional Cartan-type modular Lie superalgebras are determined, which is the outer-derivation question for the modular super case~\cite{WlTheOuterDerivationAlgebrasOfFi2020}.

The irreducible representations of Lie superalgebras of vector fields and related algebras are described~\cite{WlIrreducibleRepresentationsOfLi2020}.

The second cohomology space of the orthosymplectic Lie superalgebra is computed~\cite{WlSecondCohomologySpaceOfTheOrth2020}.
